\basicheader
\usepackage{ifthen,etoolbox}
\usepackage{colortbl}
\usepackage{longtable}
\usepackage{luacode}
\usepackage{xstring}

\begin{luacode}
    mx = 0
    my = 0
    r = 0
    delta = .42
    textwidth = 0
    off = 35

    function get(s)
        return tonumber(string.sub(s, 1, string.len(s) - 2))
    end

    function setwidth(s)
        textwidth = get(s)
    end

    function compute(_x, _y)
        x = get(_x)
        y = get(_y)

        mx = x / 2
        my = 0
        r = x / 2
    end

    function rc(ynorm)
        tex.print("(")
        x = math.sqrt(1 - ynorm * ynorm)
        tex.print(mx + x * r + 8.5)
        tex.print("pt, ")
        tex.print(my + ynorm * r)
        tex.print("pt)")
    end

    function lc(ynorm)
        tex.print("(")
        x = math.sqrt(1 - (ynorm  +delta/2) * (ynorm + delta/2))
        print(x)
        tex.print(-mx - x * r - 11)
        tex.print("pt, ")
        tex.print(my + ynorm * r)
        tex.print("pt)")
    end

    function rmid()
        rc(-.16)
    end

    function rtop()
        rc(-.16 + delta)
    end

    function rbot()
        rc(-.16 - delta)
    end

    function rrc(ynorm)
        tex.print("(")
        tex.print(textwidth - off)
        tex.print("pt, ")
        tex.print(ynorm * r + my)
        tex.print("pt)")
    end

    function rrtop()
        rrc(-.16)
    end

    function rrmid()
        rrc(-.16 - delta)
    end

    function rrbot()
        rrc(-.16 - 2* delta)
    end
\end{luacode}
\AtBeginDocument{\directlua{setwidth("\the\linewidth")}}

\setlength{\LTpost}{\medskipamount}

\usepackage{enumitem}
\ifthenelse{\equal{\TemplateLanguage}{he}}{
\AtBeginDocument{\def\footnoterule{\vskip.25cm\hbox to \textwidth{\kern.6\textwidth\hrulefill}\vskip.25cm}}
\setlist[itemize]{leftmargin=0pt, rightmargin=1cm}
\setlist[description]{leftmargin=0pt, rightmargin=0pt}
}{
\usepackage[sfdefault,tabular,scaled=.9]{FiraSans}
}
\setlist[itemize,1]{label={\begin{tikzpicture}
    \useasboundingbox (0,-.5pt) rectangle (5pt,4.25pt);
    \fill[myyellow] (0,0)--(0,5pt)--(5pt,2.5pt)--cycle;
\end{tikzpicture}}}


\usepackage[lining,scaled=.9]{FiraMono}
\usepackage[mathrm=sym]{unicode-math}
\usepackage[Scale=.9]{firamath-otf}
\usepackage{luatexbase}
\usepackage{microtype}
\usepackage[hang,symbol,multiple]{footmisc}
\usepackage{tikz}
\usepackage{siunitx}
\usepackage{xcolor}

\renewcommand{\footnotemargin}{1em}
\addtolength{\textheight}{2.5ex}

\input graphdrawing.tex

\input general/protected.tex
\input general/translation-\TemplateLanguage.tex

\ifthenelse{\equal{\TemplateLanguage}{da}}{
    \usepackage[danish]{babel}
}{
\ifthenelse{\equal{\TemplateLanguage}{de}}{
    \usepackage[ngerman]{babel}
}{
\ifthenelse{\equal{\TemplateLanguage}{en}}{
    \usepackage[english]{babel}
}{
\ifthenelse{\equal{\TemplateLanguage}{et}}{
    \usepackage[estonian]{babel}
}{
\ifthenelse{\equal{\TemplateLanguage}{fi}}{
    \usepackage[finnish]{babel}
}{
\ifthenelse{\equal{\TemplateLanguage}{he}}{
\usepackage[nil,bidi=basic]{babel}
\addtokomafont{disposition}{\rmfamily}
\setmainfont{FiraGO}[Numbers=Lowercase, Scale=.9]
\setsansfont{FiraGO}[Numbers=Lowercase, Scale=.9]
\babelfont{rm}[Numbers=Lowercase, Scale=.9]{FiraGO}
\babelfont[hebrew]{rm}[Numbers=Lowercase, Scale=.9]{FiraGO}
}{
\ifthenelse{\equal{\TemplateLanguage}{nb}}{
    \usepackage[norsk]{babel}
}{
\ifthenelse{\equal{\TemplateLanguage}{pl}}{
    \usepackage[polish]{babel}
}{
\ifthenelse{\equal{\TemplateLanguage}{sv}}{
    \usepackage[swedish]{babel}
}{
\ifthenelse{\equal{\TemplateLanguage}{uk}}{
    \usepackage[ukrainian]{babel}
}{}}}}}}}}}}

\renewcommand{\setlang}[1]{\ifthenelse{\equal{\TemplateLanguage}{he}}{\begin{otherlanguage}{hebrew}#1\end{otherlanguage}}{#1}}
\ifthenelse{\equal{\TemplateLanguage}{he}}{
\rhead{\setlang{\hskip-19pt
\directlua{setwidth("\the\linewidth")}
\begin{tikzpicture}
    \node[inner sep=0pt] (A) at (0,0) {\includegraphics[scale=.475]{header.png}};
    \useasboundingbox (A.north west) -- (A.north east) -- (A.south east) -- (A.south west) -- cycle;
    \path let \p1 = (A.north east) in \directlua{compute("\x1", "\y1")} {};
    \node[anchor=south west, inner sep=0pt] (B) at \directlua{rtop()} {\setlang{\textbf{\fontsize{22pt}{30pt}\newfontface\titling{Fira Sans Light}[Numbers=Uppercase]\titling\selectfont$\!$CEOI \olympiadyear}\vphantom{lg}}};
    \node[anchor=south west, inner sep=0pt] (C) at \directlua{rmid()} {\setlang{\tCEOIPlace\vphantom{lg}}};
    \node[anchor=south west, inner sep=0pt] at \directlua{rbot()} {\setlang{\tCEOIDates\vphantom{lg}}};

    \node[anchor=south east, inner sep=0pt] at \directlua{rrtop()} {\setlang{\textbf{\contestname}\vphantom{lg}}};
    \node[anchor=south east, inner sep=0pt] at \directlua{rrmid()} {\setlang{{\tTask}:\enspace\textbf{\taskname}\vphantom{lg}}};
    \node[anchor=south east, inner sep=0pt] at \directlua{rrbot()} {\setlang{{\tLang}:\enspace\textbf{\TemplateLanguage}\vphantom{lg}}};
\end{tikzpicture}}}
}{\lhead{\hskip-17pt
\directlua{setwidth("\the\linewidth")}
\begin{tikzpicture}
    \node[inner sep=0pt] (A) at (0,0) {\includegraphics[scale=.475]{header.png}};
    \useasboundingbox (A.north west) -- (A.north east) -- (A.south east) -- (A.south west) -- cycle;
    \path let \p1 = (A.north east) in \directlua{compute("\x1", "\y1")} {};
    \node[anchor=south west, inner sep=0pt] (B) at \directlua{rtop()} {\setlang{\textbf{\fontsize{22pt}{30pt}\newfontface\titling{Fira Sans Light}[Numbers=Uppercase]\titling\selectfont$\!$CEOI \olympiadyear}\vphantom{lg}}};
    \node[anchor=south west, inner sep=0pt] (C) at \directlua{rmid()} {\setlang{\tCEOIPlace\vphantom{lg}}};
    \node[anchor=south west, inner sep=0pt] at \directlua{rbot()} {\setlang{\tCEOIDates\vphantom{lg}}};

    \node[anchor=south east, inner sep=0pt] at \directlua{rrtop()} {\setlang{\textbf{\contestname}\vphantom{lg}}};
    \node[anchor=south east, inner sep=0pt] at \directlua{rrmid()} {\setlang{{\tTask}:\enspace\textbf{\taskname}\vphantom{lg}}};
    \node[anchor=south east, inner sep=0pt] at \directlua{rrbot()} {\setlang{{\tLang}:\enspace\textbf{\TemplateLanguage}\vphantom{lg}}};
\end{tikzpicture}}}

\addtolength{\footskip}{-1.5ex}
\cfoot{{\textbf{\thepage}}/\pageref{LastPage}}

\addtolength{\topmargin}{-5pt}

\makeatletter
\def\sthelper#1#2{{\sffamily\bfseries \tSubtask{#1} (\IfInteger{#2}{\ifnum#2=1\tPoint{#2}\else\tPoints{#2}\fi}{\tPoints{#2}}).\enspace\hss}}
\newcount\stcount
\stcount=0
\def\st#1{\setbox0=\hbox{\sthelper{9}{99}}\global\advance\stcount by 1 \par\ifdim\lastskip<\smallskipamount\removelastskip\smallskip\fi\noindent\hangindent=\wd0 \hangafter=1%
\setbox1=\hbox{\sthelper{\the\stcount}{#1}}%
\ifdim\wd1>\wd0\box1\else\hbox to \wd0{\sthelper{\the\stcount}{#1}}\fi\ignorespaces}
\def\subtask{\count1=\stcount \advance\count1 by 1 \st{\subtaskpoints{\the\count1}}}
\def\currconstraint#1{\scopedconstraint{\the\stcount}{#1}}
\def\currconstraints{\@ifstar{\currconstraint{@ll*}}{\currconstraint{@ll}}}
\def\currconstraintupper#1{\scopedconstraintupper{\the\stcount}{#1}}
\def\currconstraintlower#1{\scopedconstraintlower{\the\stcount}{#1}}
\def\currconstraintvalue#1{\scopedconstraintvalue{\the\stcount}{#1}}

\def\programcolumnwidth{3.5cm}
\def\returncolumnwidth{3cm}
\def\explanationcolumnwidth{7.55cm}
\begingroup
\catcode`\ǁ=4
\catcode`\& =\active%
\gdef\@ctivateAMP{\def&{ǁǁ}}
\endgroup
\newenvironment{interactiontable}{\begingroup%
\catcode`\& =\active%
\@ctivateAMP%
\par\medskip%
\ifthenelse{\equal{\TemplateLanguage}{he}}{\begin{otherlanguage}{english}}{}
\begin{tabular}{>{\cellcolor[gray]{.9}}p{\programcolumnwidth}@{\hskip0pt}p{.52cm}@{\hskip0pt}>{\cellcolor[gray]{.9}}p{\returncolumnwidth}@{\hskip0pt}p{.52cm}@{\hskip0pt}>{\cellcolor[gray]{.9}\sffamily\hangindent=1.25em\hangafter=1\raggedright\arraybackslash}p{\explanationcolumnwidth}}
\hline
\multicolumn{1}{|c}{\cellcolor{white}\sffamily\tProgram}&&\multicolumn{1}{c}{\sffamily \tReturn}&&\multicolumn{1}{c|}{\sffamily\tExplanation}\\\hline
\noalign{\smallskip}%
}{\end{tabular}
\ifthenelse{\equal{\TemplateLanguage}{he}}{\end{otherlanguage}}{}\endgroup\par\medskip}


\def\myc@ption#1#2#3{\noalign{\bfseries\large #3}}
\def\showcases{\removelastskip\par\bigskip\begingroup%
% This is evil: we hijack LT@makecaption to force longtable to place the heading
% same page as the table...
\let\LT@makecaption=\myc@ption
\begingroup%
\microtypesetup{activate=false}%
\testcasetable%
\endgroup%
\endgroup}
\def\showlimits{\sheading{\tLimits}%
\tTime: \timelimit\\
\tMemory: \memlimit}
% ---
\makeatother

\def\flushsubtasks{\global\stcount=\numsubtasks}

\AtEndDocument{\ifnum\stcount=\numsubtasks\else\PackageError{lg-template}{Not all subtasks are contained in the task statement!\MessageBreak You should call \protect\subtask\space (or \protect\st) precisely as often as there are subtasks, not counting the public one. If for some reason you do not want to do this, you can call \protect\flushsubtasks\space after all calls to \protect\subtask\space and \protect\st}{You should call \protect\subtask\space (or \protect\st) precisely as often as there are subtasks, not counting the public one. If for some reason you do not want to do this, you can call \protect\flushsubtasks\space after all calls to \protect\subtask\space and \protect\st}\fi}

\definecolor{myred}{RGB}{230, 37, 46}
\definecolor{mygreen}{RGB}{46, 168, 83}
\definecolor{myyellow}{RGB}{249, 187, 32}
\definecolor{myblue}{RGB}{46, 83, 168}

\def\setmymathdims{\global\Umathsubshiftdown\textstyle=2.5pt%
\global\Umathsubshiftdrop\textstyle=2.5pt%
\global\Umathsubtopmax\textstyle=4pt%
\global\Umathsupshiftup\textstyle=3.75pt%
\global\Umathopopenspacing\textstyle=225mu
\global\Umathsubshiftdrop\displaystyle=3pt%
\global\Umathsubshiftdown\displaystyle=3pt%
\global\Umathsubtopmax\displaystyle=3pt%
\global\Umathsupshiftdrop\displaystyle=4pt%
\global\Umathsupshiftup\displaystyle=4pt%
\global\Umathopopenspacing\textstyle=2.25mu}
\AtBeginDocument{\everymath{\setmymathdims}\setbox0=\hbox{$ $}\setmymathdims}